\section{原创创新点拆解}

本研究从《海洋大数据技术原理与应用》课程核心议题出发，经历了问题提出、实证证伪与问题锐化的完整演进。课上讲授的“海洋大数据的计算引擎”使作者联想到：自主水下航行器端侧的算力与能量十分有限，计算、测量与运动之间必须权衡。作者最初提出的问题是：当航行器带着几天前的旧卫星图下水时，应把时间花在继续机载推演上，还是停下来重新观测？在东海赤潮巡测前期研究（第3节）中，我们发现增加推演计算只会让规划器在想象模型里自评虚高，真实收益早已饱和，且平滑标量场中用时分配的改进空间仅约0.02，远小于下水位置与地图新鲜度的影响。针对前期方法增益小、新意不足的缺陷，作者将问题锐化并转向强海流绕障这一选错一侧就损失数小时的二元决策：当存在多个与旧图同样相容、却要求采取相反行动的海洋状态时，再多计算也无法区分；能否在历史大数据界定的合理流场范围内，专门搜出一个让另一侧更好、且使当前选择明显受损的合理反例？这一追问直接催生了本文的合理反例检验框架（表\ref{tab:innovation}）。

\begin{table}[htbp]
\centering
\caption{本文核心原创创新点、已有做法对比、实验证据与局限拆解。}
\label{tab:innovation}
\footnotesize
\begin{tabular}{@{}p{0.13\linewidth}p{0.20\linewidth}p{0.22\linewidth}p{0.19\linewidth}p{0.14\linewidth}@{}}
\toprule
创新点 & 已有做法及其不足 & 本文做法 & 证据（图表） & 局限 \\
\midrule
1. 分离数值歧义与环境歧义的决策判据 & 确定性规划在粗图上盲目加密网格（误以为多算能消除一切模糊）；传统元推理只管何时停止搜索，不区分误差源于求解器还是环境模型失真。 & 提出数值歧义（小于粗精网格容限 $\tau_{\text{num}}$，靠按需精算消除）与环境歧义（合理集合内存在翻转决策的反例，只能靠新信息或保守应对）的二分准则。 & 图\ref{fig:overview}、图\ref{fig:curse}、图\ref{fig:ambiguity}(a)：64倍精算仅将选错率从21.7\%降至19.7\%，当天0.5°图则降至10.4\%。 & 仅在左右绕障的二元拓扑决策上验证，尚未扩展至多分支同伦路径。 \\
\addlinespace
2. 基于包络定理梯度的合理反例主动搜索 & 随机合理场按期望时间选边需上千个样本（约2048次求解）才能压住抽样噪声，32个样本时改选多为噪声；翻转比例只计翻转多少而不计损失多大。 & 利用包络定理沿当前最快航线累加流速导数，一次梯度仅需2次最短路求解；先沿初始梯度跳至球面再投影细化，直接求最坏后悔。 & 图\ref{fig:risk}(a)、图\ref{fig:budget}：约71次求解即达到1024个随机合理场（2048次求解）的选边效果；最坏后悔与期望后悔排前20\%的任务均捕获约三分之二的损失，明显高于翻转比例（44\%--49\%）。 & 随机合理场的期望后悔已能较好地排序风险任务，故反例搜索的优势集中在减少选边所需的求解次数；大扰动下最短路拓扑可能跳变，局部梯度近似会产生折耗。 \\
\addlinespace
3. 海洋大数据驱动的合理集合与半径标定 & 若凭经验设定不确定性集合的形状与半径，会脱离海洋物理结构，易过度保守或漏报黑潮切变。 & 用三年历史再分析误差标定各点误差大小（缩放矩阵 $\mathbf{S}(\boldsymbol{x})$）与76千米相关长度的无散流函数模态；以覆盖90\%以上真实两侧时间差标定最小半径 $r=1.5$。 & 图\ref{fig:data}、图\ref{fig:ambiguity}(b)、表\ref{tab:scope}：$r=1.5$在测试年份覆盖94\%真实差值；消融替换为历史相似日差值后保守增益消失（0.99 vs 0.95小时）。 & 依赖多年历史再分析数据标定平稳统计量，对台风过境等非平稳极端天气适应受限。 \\
\addlinespace
4. 权衡最坏后悔与观测代价的门控响应 & 自适应采样面向全场降方差去测流，或不论代价一律保守/一律上浮，忽视了原位测流看不见远方切变及上浮通信本身耗时。 & 以最坏后悔超阈值 $\theta=8$小时作为门控：超标且上浮代价低（$<1$小时）时上浮取新图，代价高或仅有原位测流时改选最坏后悔更小一侧。 & 图\ref{fig:main}、图\ref{fig:risk}(b)、表\ref{tab:scope}：黑潮区门控后取保守将平均损失从0.98降至0.70小时，上浮代价切换点约1小时；湾流区与阿古拉斯回流区门控后上浮取图分别将损失从3.56降至1.31小时、从6.87降至2.20小时。 & 不观测的保守改选在湾流与阿古拉斯回流区这类两侧都可能被强流堵死的海区均无显著收益，此时须靠门控触发观测；黑潮区改带真实卫星图后其收益也不再显著。 \\
\bottomrule
\end{tabular}
\end{table}

在人机协同完成本研究的过程中，作者与人工智能工具的分工如下（详见附录B，对话记录见代码仓库）。作者提出了最初关于机载计算与观测如何权衡的问题；在第一版结果出来后要求改用真实卫星数据检验、补充对照与消融实验；在前期研究结果出来后指出用时分配增益过小、方法新意不足，并提出了“两条航路、一张粗图、能否找到让另一条路更好的合理海流反例”这一新问题，决定转向本文的研究。人工智能助手Claude（Anthropic）完成了文献检索、代码、实验与作图，并核对了正文中的数字与引用；Gemini（Google）根据Claude整理的事实清单与提纲撰写了论文文字。
